\section{Visualization}
\label{sec:visualization}

This section
presents the key visualizations produced by our analysis framework, illustrating the qualitative behavior of the models and
experiments described above.

\subsection{Model Comparison}

Figure~\ref{fig:metric_chart} summarizes the precision-recall trade-off across all three models on the test segment. The GP Winner favors precision over recall; our UNet2D achieves higher recall at comparable precision; our TimeSformer reaches the highest precision but the lowest recall among the 100-epoch models.

\begin{figure}[h]
\centering
\includegraphics[width=0.85\textwidth]{metric_comparison_chart}
\caption{Metric comparison across models (test segment 20230504223647). The UNet2D achieves the most balanced precision-recall trade-off among the 100-epoch models.}
\label{fig:metric_chart}
\end{figure}

The difference map (Figure~\ref{fig:diff_map}) shows where models disagree spatially. Positive differences (red) concentrate along character edges and faint traces, reflecting our model's higher recall. Negative differences (blue) are sparse, appearing where the GP's 26-layer context provides additional discriminative power not captured by our 5-layer window.

\begin{figure}[h]
\centering
\includegraphics[width=0.65\textwidth]{diff_map}
\caption{Difference map (UNet2D minus GP Winner) on the test segment. Red = our model detects ink the GP misses; blue = reverse. Positive differences dominate, consistent with higher recall.}
\label{fig:diff_map}
\end{figure}

\subsection{Depth Window Position Effect}

Figure~\ref{fig:sweep_comparison} shows predictions at three Z-window positions on the test segment. At the optimal $s = 30$, predictions are densest and most coherent. Shallower windows introduce pre-surface noise; deeper windows rapidly lose ink signal as the window moves below the papyrus surface consistent with the quantitative results in Table~\ref{tab:exp_b}.

\begin{figure}[h]
\centering
\begin{tabular}{ccc}
\includegraphics[width=0.30\textwidth]{sweep_s28_test} &
\includegraphics[width=0.30\textwidth]{sweep_s30_test} &
\includegraphics[width=0.30\textwidth]{sweep_s32_test} \\
(a) $s = 28$ (layers 28--32) & (b) $s = 30$ (layers 30--34) & (c) $s = 32$ (layers 32--36) \\
\end{tabular}
\caption{Position sweep on the test segment. Model weights are identical; only the input window shifts. Window~(b) is the optimal configuration.}
\label{fig:sweep_comparison}
\end{figure}

\subsection{Leave-One-Out Sensitivity}

Figure~\ref{fig:loo_visual} contrasts full-layer predictions against predictions with layer~32 masked. Confidence drops selectively in ink-bearing regions while background areas remain largely stable, confirming that layer~32's contribution is signal-driven rather than incidental.

\begin{figure}[h]
\centering
\begin{tabular}{cc}
\includegraphics[width=0.45\textwidth]{loo_baseline_train} &
\includegraphics[width=0.45\textwidth]{loo_mask_ch2_train} \\
(a) Baseline (all 5 layers) & (b) Layer 32 masked \\
\end{tabular}
\caption{Leave-one-out visualization (training segment 20230504231922). Masking layer~32 reduces prediction confidence in ink regions while background predictions remain stable.}
\label{fig:loo_visual}
\end{figure}

\subsection{Architecture Comparison}

Figure~\ref{fig:arch_comparison} compares all three architectures on the test segment. The UNet2D produces the broadest coverage; both TimeSformer variants show sparser, higher-precision predictions. With only 5 input layers, channel-wise convolution is well-matched to this shallow depth stack, while the attention advantage of TimeSformer requires more frames to be effective.

\begin{figure}[h]
\centering
\begin{tabular}{ccc}
\includegraphics[width=0.30\textwidth]{pred_gp_winner} &
\includegraphics[width=0.30\textwidth]{pred_unet2d_100ep} &
\includegraphics[width=0.30\textwidth]{pred_timesformer_100ep} \\
(a) GP Winner & (b) Our UNet2D (100ep) & (c) Our TimeSformer (100ep) \\
\end{tabular}
\caption{Architecture comparison (test segment). (a)~GP Winner ($F_{0.5} = 0.379$, 26 layers). (b)~Our UNet2D ($F_{0.5} = 0.341$, 5 layers). (c)~Our TimeSformer ($F_{0.5} = 0.290$, 5 layers). The UNet2D provides the widest coverage; both TimeSformer variants favor precision over recall.}
\label{fig:arch_comparison}
\end{figure}
